\documentclass[10pt,openany]{book}
\input{../preamble}
\renewcommand{\thesection}{\arabic{section}}
\renewcommand{\thesubsection}{\thesection.\arabic{subsection}}
\fancyhead[L]{\scriptsize\sffamily はじめての銀河シミュレーション}
\fancyhead[R]{\scriptsize\sffamily 構成案 v9}
\begin{document}
\begin{center}
{\Large\sffamily\bfseries\color{navy} Pythonによる\\はじめての銀河シミュレーション}\par\medskip
{\large テスト粒子で学ぶ重力・軌道・銀河相互作用}\par\bigskip
{\Large\sffamily 構成案 Version 9}\par\medskip
参考図書整理版 v2.5 / 2026年10月8日
\end{center}
\section{本書のねらい}
理学部の学部生が、銀河力学・天体力学のゼミで学んだ重力、軌道、潮汐相互作用を、Pythonによる計算へ移すための実習型教材です。第0章でPython環境を準備し、Jupyter Notebookの利用法を確認した後、1粒子、2体問題、銀河中の軌道、円盤、銀河相互作用へ段階的に進みます。

最終モデルは、二つの移動するPlummerポテンシャルと質量を持たないテスト粒子です。円盤の自己重力、ガス、星形成、宇宙論的初期条件は扱いません。書名から「形成」を外し、計算する内容との対応を明確にします。基本計算は共面2次元で、円盤傾斜角や鉛直構造の比較も範囲外です。

数値計算法は、問題に必要な場所で主な方法だけを説明します。Euler法は導入、KDK leapfrogは重力計算の標準、RK4は一般ODEと速度依存力のための主要手法です。台形則、二分法、線形補間と逆変換標本化も、銀河モデルの中で使います。

\section{全体構成}
\begin{center}\small
\begin{tabularx}{\linewidth}{cXX}\toprule
章 & 主題 & その章で完成するもの\\\midrule
0 & 計算環境を準備する & Python・Jupyterの準備\\
1 & Pythonで天体の運動を計算する & 1粒子のKDK・RK4積分\\
2 & 2体問題を計算する & 重心・相対運動と2体軌道\\
3 & 銀河の中で星を動かす & 質量積分・求根・Plummer軌道\\
4 & テスト粒子で円盤銀河を作る & CDFから回転円盤を生成\\
5 & 2つの銀河を接近させる & 潮汐テイル・ブリッジ\\
6 & 銀河相互作用を調べる & 条件を変えた比較実験\\
7 & 動力学的摩擦で軌道減衰を調べる & 速度依存力による衛星軌道の減衰\\\bottomrule
\end{tabularx}\end{center}
各章は「問い → モデル → 式 → 計算手順 → 実装 → 実験 → 解釈と確認」で進みます。章・節・小節の見出しは、原則として学生が行う操作や確認が分かる動詞句とし、「本章のまとめ」「演習」「参考図書」は定型見出しとして残します。本文、完成版Notebook、学生版Notebookの役割を分け、章ごとに独立実行できる形を保ちます。

\clearpage
\section{第0章　計算環境を準備する}
第1章以降の学習を始める前の準備だけを扱います。Python 3、NumPy、Matplotlib、Jupyterを利用できる環境を用意し、配布された\texttt{requirements.txt}から必要なパッケージを導入します。仮想環境の作成方法と、macOS・Linux・Windowsでの有効化方法を簡潔に示します。

JupyterLab / Jupyter Notebookの起動、Codeセルの実行、カーネルの再起動、全セルの再実行、完成版・学生版Notebookの役割を説明します。最後に版番号を表示する短いセルで動作確認を行います。Python文法、NumPy配列、関数、乱数、Matplotlib、浮動小数点数の一般的な解説は置きません。本章にはまとめ・演習を設けません。

\section{第1章　Pythonで天体の運動を計算する}
\textbf{学習目標：}Newton方程式を1階の連立方程式へ直し、KDKとRK4を実装できます。

固定中心質量中の1粒子を使い、Euler法からKDKへ進みます。半kick、drift、半kickの更新と加速度再利用を説明します。RK4も4段階の予測状態から導き、Kepler円軌道で比較します。

\textbf{確認：}同時刻の解析位置、エネルギーと角運動量、時間刻み、長時間の誤差、力評価回数を使います。同じ刻みと同じ力評価回数の比較を分けます。TikZでKDKの操作順序とRK4の予測状態を示します。

\section{第2章　2体問題を計算する}
\textbf{学習目標：}二つの天体を作用・反作用を保って動かし、重心と相対軌道を区別できます。

質量比2:1の点質量から始め、円・楕円・放物線を作ります。第1章のKDKを再利用し、近点を指定した内向きの接近軌道も作ります。新しい時間積分法は追加しません。

\textbf{確認：}全エネルギー、全角運動量、全運動量、重心、解析的な円軌道と楕円周期を比較します。初期エネルギーが0の放物線では別の正の尺度で規格化します。TikZで原点・重心・相対位置と、力・加速度の違いを示します。

\clearpage
\section{第3章　銀河の中で星を動かす}
\textbf{学習目標：}Plummerモデルを用い、台形則・二分法・時間積分を別の目的で使えます。

密度から内側の質量を積分し、二分法で半質量半径を求めます。Plummerの解析式が分かる例で、数値積分と求根を検証します。後半では解析的加速度をKDKへ渡し、円軌道と非円軌道を動かします。

\textbf{確認：}積分範囲、分点数、求根許容幅、時間刻みを別々に変えます。小さな求根残差と真の半径誤差を区別します。球殻・二分法の図と、台形の高さ・幅の注記を置きます。

\section{第4章　テスト粒子で円盤銀河を作る}
\textbf{学習目標：}指数円盤の数面密度を確率密度へ変え、線形補間と乱数から多数粒子を配置できます。

台形則でCDFを作り、線形補間を数式と1点の実装で学び、\texttt{np.interp}による逆変換へ進みます。半径と方位角に独立な乱数を使い、Plummer場の円運動速度を与えます。加速度を多数粒子へ拡張し、KDKを再利用します。

\textbf{確認：}CDF表の近似、有限標本、時間積分の誤差を区別します。数面密度、粒子の保存量、円運動の解析位置を調べます。円環の面積と逆補間の図を使います。粒子分布と背景の質量分布を混同しません。

\section{第5章　2つの銀河を接近させる}
\textbf{学習目標：}第2--4章の部品を統合し、動く中心と星を同じ時刻へ進められます。

中心間の有効対ポテンシャルを定義し、ゼロエネルギー・指定近点の初期軌道を作ります。二つの局所円盤に中心の位置・速度を足し、全行をKDKで同期更新します。標準は等質量、$a_1=a_2=b=3$ kpc、$D_0=60$ kpc、$D_{\rm p}=8$ kpc、各1000粒子です。

\textbf{確認と解釈：}中心二体系の保存量、星の自己収束、出身中心から外へ移った割合を調べます。時間依存する場の中の星にはエネルギー保存を要求しません。標準は円盤の広がりが重なる強い遭遇で、中心間の有効モデルは厳密な拡がった球同士の相互力ではありません。力の関係、位置・速度の並進、回転の向き、潮汐加速度の差をTikZで示します。

\clearpage
\section{第6章　銀河相互作用を調べる}
\textbf{学習目標：}変更した入力、連動する初期値、測定条件を整理し、小さな比較実験を説明できます。

回転方向、近点距離、質量比を独立な系列として比較します。近点通過からの時間をそろえ、外側粒子割合を定義して測ります。閾値、seed、時間刻みを変え、差の由来を確かめます。非単調な近点系列も残し、中心間の有効長さによる力の比を追加の確認として示します。その比だけで非単調性の原因は決めません。

\textbf{確認：}同じ物理モデルの中で、物理パラメータ、有限標本、測定条件、時間刻みを区別します。TikZで近点からの比較時刻のそろえ方を示します。本章の最終課題は、三つの系列から一つを選び、小さな研究として報告することです。

\section{第7章　動力学的摩擦で軌道減衰を調べる}
\textbf{学習目標：}新しい物理として速度依存の動力学的摩擦を加え、RK4で軌道減衰を計算できます。

固定Plummer主銀河と質量比1:10の衛星へモデルを切り替え、Chandrasekhar動力学的摩擦を衛星中心に加えます。背景密度と速度分散、Coulomb対数、局所Maxwell近似を区別し、摩擦あり・なしの軌道、エネルギーと仕事、角運動量、時間刻みを比較します。追加実習では衛星のテスト粒子も動かしますが、星には摩擦を直接与えません。

\textbf{確認：}数値解の収束と摩擦モデルの妥当性を分け、背景応答、衛星質量損失、主銀河の変形を省くことを明示します。TikZでは、直接計算していない背景応答と、実装する近似摩擦項を区別します。

\section{巻末資料}
全章の演習解答は巻末の「演習の解答・着眼点」にまとめ、参考図書も一か所へ集約します。銀河力学はRydenとBinney \& Tremaineを基本とします。数値計算法は杉浦、伊理・藤野、『シミュレーション天文学』を主に参照し、Pythonでの実装例を補う参考書として幸谷『Python数値計算プログラミング』を挙げます。

\section{教材の制作と更新}
コードはJupytext master、文章はLaTeXを正本とします。本文のコード・実測値・派生値は生成物から読み込み、完成版と学生版Notebookも同じmasterから作ります。第1--7章ではgoalbox、図表・式番号を統一し、各章末を「本章のまとめ→演習」とします。第0章は準備章なのでgoalbox・まとめ・演習を置きません。演習解答と参考図書は巻末へ集約します。具体的な規約はスタイルガイドv1.8に従います。

標準モデルの重力・初期条件・刻みは全章改訂前を維持します。第5章は0.05 Myr、0.60 Gyr、12000ステップを基準とし、パラメータや積分時間を変更した場合は、その条件で刻みを確認します。第7章の中心のみの摩擦例は1 Myr、1.50 Gyrです。これらを任意の計算に対する精度保証としません。

最終課題は、問い、モデルと条件、数値設定、結果、独立な確認、解釈と限界の順にまとめます。静止画を少なくとも一つ、定義を示した定量指標を少なくとも一つ使います。結論をその実験の条件とモデルの範囲に限定することが、本書の到達点です。
\end{document}
